%! Principal Component Analysis — Unit 7, Lesson 7.3 — Group Activity
\documentclass[10pt]{article}
\usepackage{linalg-article}
\usepackage{linalg-boxes}
\newcommand{\TallMath}[1]{$\displaystyle #1\rule[-1.4em]{0pt}{3.2em}$}
\newcommand{\uu}{\mathbf{u}}
\newcommand{\vv}{\mathbf{v}}
\newcommand{\ee}{\mathbf{e}}
\newcommand{\zero}{\mathbf{0}}
\newcommand{\T}{^{\top}}

\begin{document}

\pageheader{Unit 7, Lesson 7.3}{Group Activity}
\namepartnerperiod

\vspace{0.10in}

\begin{headlinebox}{blush}
\small\textbf{Finding the principal components.} PCA on a data table: \textbf{center} it (subtract each
column's mean), form the symmetric $A\T A$, and take its eigenvectors --- the \textbf{principal components}
$\vv_1\perp\vv_2$, biggest $\sigma$ first. The eigenvalue $\sigma^2$ is the spread along each direction. Work
with your partner, show every step, and \emph{justify} each answer.
\end{headlinebox}

\vspace{0.08in}

\begin{tcolorbox}[colback=white, colframe=black!40, title=\textbf{Tier R --- Remediate},
    fonttitle=\bfseries, arc=1.5mm, left=3mm, right=3mm, top=2mm, bottom=2mm, breakable]
Warm up the two set-up moves: center a data table, then form $A\T A$.
\begin{enumerate}[leftmargin=1.7em, itemsep=9pt]
  \item \textbf{Center.} Four data points are $(5,7),(7,5),(3,1),(1,3)$. Find each column's mean, then subtract to get the centered points. \writeline
  \item \textbf{Form $A\T A$.} Stack the centered points as the rows of $A$ and compute $A\T A=\blank{2.6cm}$. (You will meet this matrix again in Tier~A.) \writeline
\end{enumerate}
\end{tcolorbox}

\begin{tcolorbox}[colback=white, colframe=black!40, title=\textbf{Tier A --- Approaching Proficiency},
    fonttitle=\bfseries, arc=1.5mm, left=3mm, right=3mm, top=2mm, bottom=2mm, breakable]
Run PCA on the data whose centered points are $(1,3),(3,1),(-1,-3),(-3,-1)$, giving
$A\T A=\begin{bsmallmatrix} 20 & 12\\ 12 & 20 \end{bsmallmatrix}$.
\begin{enumerate}[leftmargin=1.7em, itemsep=9pt]
  \item \textbf{Eigenvalues.} Solve $\det(A\T A-\lambda I)=0$ to find $\lambda_1,\lambda_2$, then $\sigma_1,\sigma_2$. \writeline
  \item \textbf{Principal components.} Solve $(A\T A-\lambda I)\vv=\zero$ for each $\lambda$ to find PC1 and PC2 (as unit vectors). \writelines{2}
  \item \textbf{How much spread on PC1? (justify)} Compute the fraction $\sigma_1^2/(\sigma_1^2+\sigma_2^2)$. Would summarizing each point by its PC1 coordinate alone be a good idea here --- why or why not? \writeline
\end{enumerate}
\end{tcolorbox}

\begin{tcolorbox}[colback=white, colframe=black!40, title=\textbf{Tier E --- Extension},
    fonttitle=\bfseries, arc=1.5mm, left=3mm, right=3mm, top=2mm, bottom=2mm, breakable]
\textbf{Choosing how many components to keep.} A survey of $50$ people records $4$ features per person. After
centering and forming $A\T A$, the four eigenvalues (the spread along each principal component) are
$\sigma^2=60,\ 25,\ 10,\ 5$.
\begin{enumerate}[leftmargin=1.7em, itemsep=9pt]
  \item \textbf{Total spread.} Add the four eigenvalues. \blank{2.6cm}
  \item \textbf{Cumulative.} What percent of the total spread do the first $1$, $2$, and $3$ components capture? \writeline
  \item \textbf{Choose \& justify.} What is the smallest number of components that captures at least $95\%$? Then explain, in one sentence each, why we \emph{center} first and why we keep the \emph{largest}-$\sigma$ components. \writelines{2}
\end{enumerate}
\end{tcolorbox}

\end{document}
